% Preamble for the slides of the course
% Numerical Methods for PDEs II: finite elements and spectral methods

\usepackage[T1]{fontenc}
\usepackage[english]{babel}
\usepackage{lmodern}
\usepackage{graphicx}
\usepackage{booktabs}
\usepackage{multirow}
\usepackage{tikz}
\usetikzlibrary{arrows.meta,calc}

% Shared notation with the lecture notes (\R, \diff, \defeq, \mat, \vec, \prodH, ...)
\subimport{../header/}{math.tex}

%%%%%%%%%%%%%%%%%%%%%%%%%%%%%%%%%%%%%%%%%%%%%%%%%%%%%%%%%%%%%%%%%%%%%%%%%%%%%%%%
% Appearance
%%%%%%%%%%%%%%%%%%%%%%%%%%%%%%%%%%%%%%%%%%%%%%%%%%%%%%%%%%%%%%%%%%%%%%%%%%%%%%%%

\usetheme{default}
\usefonttheme[onlymath]{serif}

\definecolor{mnublue}{RGB}{31,64,122}
\definecolor{mnuaccent}{RGB}{190,80,20}

\setbeamercolor{structure}{fg=mnublue}
\setbeamercolor{frametitle}{fg=mnublue}
\setbeamercolor{title}{fg=mnublue}
\setbeamercolor{footline}{fg=gray}
\setbeamerfont{footline}{size=\tiny}
\setbeamerfont{frametitle}{series=\bfseries}
\setbeamertemplate{navigation symbols}{}
\setbeamertemplate{blocks}[rounded][shadow=false]
\setbeamertemplate{itemize item}[circle]
\setbeamertemplate{itemize subitem}[triangle]
\setbeamertemplate{enumerate items}[default]

\setbeamertemplate{footline}{%
    \leavevmode%
    \begin{beamercolorbox}[wd=\paperwidth,ht=2.5ex,dp=1.2ex,leftskip=1em,rightskip=1em]{footline}%
        \usebeamerfont{footline}%
        Numerical Methods for PDEs II%
        \ifnum\value{lecture}>0
            \quad\textbar\quad Session~\insertlecturenumber: \insertlecture
        \fi
        \hfill\insertframenumber
    \end{beamercolorbox}}

%%%%%%%%%%%%%%%%%%%%%%%%%%%%%%%%%%%%%%%%%%%%%%%%%%%%%%%%%%%%%%%%%%%%%%%%%%%%%%%%
% Coloured blocks (same colour code as the lecture notes)
%%%%%%%%%%%%%%%%%%%%%%%%%%%%%%%%%%%%%%%%%%%%%%%%%%%%%%%%%%%%%%%%%%%%%%%%%%%%%%%%

% \mnunewblock{environment}{Heading}{colour}
% Usage: \begin{theorem}[Name] ... \end{theorem}
\makeatletter
\newcommand{\mnunewblock}[3]{%
    \newenvironment{#1}[1][]{%
        \setbeamercolor{block title}{fg=black,bg=#3!25}%
        \setbeamercolor{block body}{fg=black,bg=#3!6}%
        \begin{block}{#2\ifx\relax##1\relax\else\ (##1)\fi}}%
        {\end{block}}}
\makeatother

\mnunewblock{definition}{Definition}{blue}
\mnunewblock{theorem}{Theorem}{orange}
\mnunewblock{lemma}{Lemma}{orange}
\mnunewblock{corollary}{Corollary}{orange}
\mnunewblock{example}{Example}{green!70!black}
\mnunewblock{remark}{Remark}{gray}

% Main steps of a proof (details are given at the blackboard)
\newenvironment{proofsteps}[1][Proof: main steps]{%
    \setbeamercolor{block title}{fg=mnublue,bg=gray!12}%
    \setbeamercolor{block body}{fg=black,bg=white}%
    \begin{block}{#1}}%
    {\end{block}}

% Things to be developed at the blackboard
\newcommand{\board}[1]{%
    \par\smallskip
    {\color{mnuaccent}\textbf{At the board:}}~#1\par}

% Key message of a frame
\newcommand{\takeaway}[1]{%
    \par\smallskip
    {\color{mnublue}\(\blacktriangleright\)~\textbf{#1}}\par}

%%%%%%%%%%%%%%%%%%%%%%%%%%%%%%%%%%%%%%%%%%%%%%%%%%%%%%%%%%%%%%%%%%%%%%%%%%%%%%%%
% Sessions
%%%%%%%%%%%%%%%%%%%%%%%%%%%%%%%%%%%%%%%%%%%%%%%%%%%%%%%%%%%%%%%%%%%%%%%%%%%%%%%%

\newcommand{\repourl}{https://github.com/dgomezcastro/UAM-MNU/blob/main/}

% \notebook{path/in/repository}{file.ipynb}
\newcommand{\notebook}[2]{\href{\repourl #1/#2}{\texttt{#2}}}

% \session{label}{subject}{title}{notes}{notebooks}{exercises}
% Starts a new lecture (see \includeonlylecture) and prints its title frame.
\newcommand{\session}[6]{%
    \lecture{#3}{#1}%
    \section{Session \thelecture: #3}%
    \begin{frame}[plain]
        \vfill
        {\footnotesize\color{gray}\MakeUppercase{#2}\par}
        \vspace{.5em}
        {\usebeamercolor[fg]{title}\Large Session \insertlecturenumber\par}
        \vspace{.3em}
        {\usebeamercolor[fg]{title}\usebeamerfont{title}#3\par}
        \vspace{1em}
        {\color{mnublue}\rule{.6\textwidth}{.6pt}\par}
        \vspace{1em}
        \small
        \begin{tabular}{@{}ll@{}}
            Lecture notes: & #4 \\
            Notebook:      & #5 \\
            Exercises:     & #6
        \end{tabular}
        \vfill
    \end{frame}}
